% !PRINTED 372
\weberpage{372}
% Source-faithful three-stroke Cayley configuration sign.
\providecommand{\weberthreeup}{\mathord{\tikz[baseline=-0.2ex,x=0.23em,y=0.23em,line width=0.45pt]{\draw (0,0)--(-1,1) (0,0)--(0,1.08) (0,0)--(1,1);}}}
Betrachten wir das syzygetische Complextripel
\weq{8}{
\begin{aligned}
&y_3 z_4,\ y_4 z_3,\ \cdot\ \cdot\ \cdot\ \cdot\\
&y_3 y_4,\ z_3 z_4,\ x_3 x_4\ \cdot\ \cdot\\
&y_3 z_3,\ y_4 z_4,\ \cdot\ \cdot\ \cdot\ \cdot,
\end{aligned}
}
in dem alle Doppeltangenten vorkommen müssen, so finden wir
$[5\,8]$ und $[5\,6]$ nicht in den beiden letzten Complexen von (8),
weil
\weq{9}{
\begin{aligned}
&y_3 z_3 [5\,8] = [1\,3]\,[2\,3]\,[5\,8],\quad y_3 y_4 [5\,8] = [1\,3]\,[1\,4]\,[5\,8]\\
&y_3 z_3 [5\,6] = [1\,3]\,[2\,3]\,[5\,6],\quad y_3 y_4 [5\,6] = [1\,3]\,[1\,4]\,[5\,6]
\end{aligned}
}
das Zeichen $\vee\mid$ haben und daher azygetisch sind. Folglich kommen
 $[5\,8]$ und $[5\,6]$ im ersten der Complexe (8) vor, und die
beiden Tripel (7) sind syzygetisch.

Endlich haben wir auch das Zeichen $\sqcap$ zu betrachten, das
wieder drei Typen giebt:
\weq{10}{
\begin{aligned}
&[1\,2]\,[2\,3]\,[3\,4] = y_2 z_3 [3\,4]\\
&[1\,2]\,[2\,3]\,[3\,8] = y_2 z_3 x_3\\
&[1\,8]\,[8\,2]\,[2\,3] = x_1 x_2 z_3.
\end{aligned}
}
Auch diese Tripel sind syzygetisch, was für die beiden letzten
unmittelbar aus dem Complextripel (3) zu ersehen ist, und für
das erste aus dem syzygetischen Complextripel
\[
\begin{aligned}
&y_2 z_3,\ x_1 x_3,\ \cdot\ \cdot\ \cdot\ \cdot\\
&y_2 x_1,\ z_3 x_3,\ \cdot\ \cdot\ \cdot\ \cdot\\
&y_2 x_3,\ z_3 x_1,\ \cdot\ \cdot\ \cdot\ \cdot,
\end{aligned}
\]
folgt, von denen die beiden letzten $[3\,4]$ nicht enthalten, weil
$y_2 x_1 [3\,4]$ und $y_2 x_3 [3\,4]$ beide das Zeichen $\vee\mid$ haben.

Hier ist aber die Ausnahmestellung der Ziffer 8 gänzlich
verschwunden, und wir kommen zu dem Resultate:

\gesperrt{Unter den Tripeln von Doppeltangenten sind die mit den
Zeichen}
\[
\mid\mid\mid,\ \sqcap
\]
\gesperrt{syzygetisch, und die mit den Zeichen}
\[
\weberthreeup,\ \triangle,\ \vee\mid
\]
\gesperrt{azygetisch.}

Hieraus erhält man sehr leicht die Zeichen für sämmtliche
syzygetische und azygetische Quadrupel:
%CONT>>
